\chapter[Sermon 24. Here Is Described the Proceeding of the Holy Spirit, and Operation, the Almighty Knowledge of God, and How the Throne of God Is Borne Up or Sustained of the Four Beasts, and What the Beasts Do.]{Here Is Described the Proceeding of the Holy Spirit, and Operation, the Almighty Knowledge of God, and How the Throne of God Is Borne Up or Sustained of the Four Beasts, and What the Beasts Do.}
\label{sermon:024}
\markright{Sermon 24. Here Is Described the Proceeding of the Holy Spirit ...}

And out of the Seat proceeded lightning and thunderings and voices, and there were seven lamps of fire, burning before the Seat, which are the seven spirits of God. And before the Seat there was a Sea of glass like unto Crystal. And in the midst of the Seat, and round about the Seat, were four Beasts full of eyes before and behind. And the first beast was like a Lion, the second beast like a calf, and the third beast had a face like a man, and the fourth beast was like a flying Eagle. And the four beasts had each one of them six wings, and round about without and within, they were full of eyes. And they had no rest day neither night, saying, “Holy, holy, holy is the Lord God Almighty, which was, and is, and is to come.”

Our Lord Jesus Christ, as the faithful pastor of his church, will reveal the destinies and wonderful calamities that will come upon the church. Therefore, to stop the mouths of those who question and are inquisitive about the judgments of God, and to persuade all to have patience in these storms of evils, he sets forth a treatise beforehand, wherein he shows that all things are done or permitted to be done by God through his most just providence, and are governed or ordered by the Lamb, with a judgment most righteous and holy. For whoever believes and remembers this, whatever circumstances he encounters, he submits himself humbly, lowly, and obediently to his God, and always cries out, “The Lord is righteous in all his ways and holy in all his works.” And this is the true state of the first part of this vision, which is presented in chapters 4 and 5. It is moreover most elegant, most pleasant, and most full of consolation. All things are more vividly set forth and perceived in such fitting and heavenly representations than they can be understood in bare words.

First, a throne is described, and indeed a heavenly throne, so that in the works, providence, and judgments of God, we should not imagine anything carnal or corrupt. Secondly, he who sits on the throne is represented to us by two colors, green and red. For God is an eternal essence belonging to all their griefs or being. The same burns in love toward humankind and wills well unto man; but to the disobedient and rebels, he is a consuming fire. And the throne is surrounded by a rainbow greener than grass, comforting us so that we should not be dismayed at the sight of that godly throne, but should always remember that he who sits in the throne, judge and governor of all, is most true and keeps his promises. And to be that same league friend. Around the throne sit twenty-four elders, who have already been signified what they are, as it were shadowed, immediately in the end of the fourth chapter, and in the fifth shall be declared what they do, or what they say. Doubtless all the saints in heaven are onlookers of the judgments and works of God. For the judgments of God are not such that they should flee the light and knowledge of saints.

Now from the throne came flashes of lightning, and thunder, and other things. In the throne are he who sits and the Lamb, that is, the Father and the Son, and from both proceeds the Holy Spirit. For by interpretation it follows immediately, which are the seven spirits of God. For the lightning, the thunderings, and other things mentioned signify, or are tokens of the Holy Spirit, which elsewhere is also said to be represented by fire, and water, and wind, and by fiery tongues. But no one will think that the Holy Spirit, which is one in substance and of the simple divine nature, should be divided into seven parts. For I told you in the first chapter how the seven spirits of God are put for the sevenfold, most full, and most perfect spirit of God.

\index{Tertullian}We have in the beginning of this vision the whole mystery of the blessed Trinity, so much as is needful for us to know, and believe, and profess. There is one Seat, in that one seat are contained the sitter, the Lamb, and the Spirit: therefore there is one divine essence and nature, and thereof is one power and majesty, one rule, because there is one throne; briefly, there is one God, true and eternal, forevermore blessed, as Moses also in the sixth chapter of Deuteronomy, and all the prophets and Apostles have everywhere taught. However, in this only and undivided substance is seen a most plain distinction of persons. For there is he that sits in the throne, and the Lamb, and from both proceeds the Holy Spirit. This mystery of the Trinity we profess in the Creed. This appears openly in the incarnation of our Lord, while the angel says to the virgin, “The Holy Spirit shall come upon thee, and the power of the Most High shall overshadow thee; and that which shall be born of thee shall be called the Son of God.” Likewise, in the baptism of Christ, a voice is heard from heaven upon the Lord: “This is my well-beloved Son.” The Holy Spirit also appears in the likeness of a dove. Whereupon the Lord commanded us also to be baptized in the name of the Father, and of the Son, and of the Holy Spirit. This profession is certain and true, and so set forth by the most manifest scriptures and lively preaching of the Apostles, like as Tertullian declares against the heretic Praxeas. We ought rather believe and cleave unto these things than to the monstrous and blasphemous Spanish sophistry of Servetus, a man most corrupt.

But especially here the entire mystery of the Holy Spirit is declared to us, and that in few words, which in the Gospel of John is expressed more fully. First, his procession is noted, which in times past was rashly affirmed to be set forth in no part of Scripture. John here says, “out of the throne proceeded lightnings, \&c.” And immediately after: “which are the seven spirits of God.” This word in Greek signifies a proceeding or going out, but John says that it proceeded or went forth. Therefore, the ancient council of Constantinople rightly decreed: that is, concerning the Holy Spirit, the Lord, that quickener, proceeds from the Father, \&c. But because the Lord himself in the Gospel, speaking of the Holy Spirit, says, “he shall glorify me: for he shall take of mine, and shall show it unto you. All things whatsoever the Father has, are mine. Therefore I said, that he shall take of mine, and shall show it unto you,” no one will understand the Spirit to proceed from the Father only, and not also from the Son, concerning which there was also a lengthy contention between the Greeks and Latins. For if he proceeds from the Father, he proceeds from the Son also. For even for the same cause at present he is shown to proceed out of the Throne. But in the Throne is not only he that sits, but the Lamb also, of whom in chapter 5 we shall add, that the Lamb has seven eyes, which are the seven spirits of God, sent into the whole world. Albeit therefore that in John 15, the Holy Spirit is said to proceed from the Father, yet it is stated beforehand: “who I (says the Son) will send unto you from my father.” To be brief, if there is one substance and nature of the Father and of the Son, I do not see how the Holy Spirit should proceed from the Father without also proceeding from the Son. Let us rather leave these scrupulous disputations to idle wits; let us believe that the Spirit proceeds from both.

Moreover, the power or effect and operation of the Holy Spirit is here also set forth and declared clearly. For first, he illuminates, when he shines upon the obedient, and terrifies the rebellious with severe threats. Secondly, he governs, when he contends against this ungodly world and rebukes it for its sins, thundering out God’s terrible judgments. Two Apostles in Mark are called the sons of thunder, or thunderers. He utters moreover wholesome voices of doctrine, exhortation, and consolation, by men, for the favor of men. Finally, where the operation of the Holy Spirit cannot sufficiently well be expressed, yet by the seventh number he encompasses and accomplishes his fullness, and says that seven fiery lamps are burning before the Seat, burning, I say, not quenched, or smoking. For the grace of the Holy Spirit is bright and full of efficacy, as was spoken of before: and where these things are found in the Throne, how should any man think that the judgments proceeding from thence should be in any part corrupt, defiled, or to be blamed? By the Holy Spirit all things are preserved, and by his providence all things are wrought.

Hereunto is added another thing, a glassy sea before the seat, in clearness and brightness representing crystal. This signifies this frail world, which is subject to God, and as it were in his sight. Also, in other places of holy scripture, because of its instability, tossing, and turmoil, it bears the figure of this variable and most unconstant world. Certainly, the state of this world is more fickle than glass. Something of this shall follow in the 15th Chapter. But whatever things are done in the world through a marvelous variety, all the same shine as in a glass before the Throne, so that God sees them all as it were in a crystal: whose eyes or knowledge the least things that be can not escape. For we shall not think that such things as are done in the world are done rashly, and by a certain fortune to hap or chance besides the knowledge of God, or to be of God unknown.

\index{Ezekiel (Book of)}After this, he returns again to the throne, intending to complete what he had begun to describe, and to show that all the works of God, done by his creatures, are most holy. Royal seats, chairs, or thrones of kings are typically borne up and adorned with beasts, as Solomon’s seat was with lions, as can be seen in the third book of Kings, chapter 10. In other places, the most excellent beasts draw the triumphant chariots of princes. In the same manner, by a figure of speech, beasts are assigned to the throne of God. For God is carried upon Cherubim in his prophets. And Ezekiel, in chapter 10, names Cherubim, beasts, and the entire text demonstrates that the passage must be understood as God’s chariot, drawn by beasts, in which he himself was carried out of the city of Jerusalem. There is much mention in poets of the chariot of the Gods, taken happily by the first writers from the holy scriptures. For Satan, the ape of God, always seeks to defame the word of truth. But we, omitting the triflings of poets, will consider the sober description of this carriage of God, or rather of God’s throne. Almighty God sits in this seat. Sitting in the Scriptures signifies government. Here it is signified that God sits in all his creatures; that is to say, governs his creatures, and by his most wise providence works all in all, using every creature according to his good and just pleasure, according to the nature of each. We shall say then that by those beasts are understood all the creatures of God, dispersed throughout the four quarters of the world, which are comprehended in the whole world.

And first, it is shown in what place of the throne the beasts were: namely, in the midst of the throne, and in the circuit of the same. You will ask, if they are in the throne, how should they be about the throne? If they are about the throne, how are they in the midst of the throne? The thing must be conceived as I admonished before, that we understand that under the throne, the midst of the beasts does with their hind parts reach to the midst of the throne inwardly, and so as it were to have borne up the throne. And with their fore parts, I mean, with their breasts, and heads and wings, to have stood forth, and so to have compassed the throne, and as it were environed it round about. For so they might seem to be in the midst of the same throne, and round about the same.

After, the manner of beasts is described diligently, and they were four in number. For in times past also the number was expressed by Ezekiel. And the parts of the world are suitably signified by the fourth number, comprehending the universality of things. Some here have suggested the four monarchies of the world, \&c. And every beast had his face and his body, six wings, and was full of eyes within, as also their bodies were full of eyes. The first was like a lion in shape and fashion, the second like a calf, the third like a man, and the fourth like a flying eagle. By these appear to be signified all creatures, visible and invisible, reasonable and unreasonable, and the most excellent. For afterward in chapter five we shall hear that all creatures jointly worship the Lamb and him that sits on the Throne. And truly, God uses them all—the sun, the moon, the stars, the air, the fire—and briefly, all living things. And such creatures as he has chosen, in order to work something by them, he makes them efficacious, instructing each according to their state and condition, that they should lack no wisdom, reason, strength, power, patience, labor, quickness, nor swiftness. The face of man signifies wit and wisdom, as also the eyes signify foresight, watchfulness, subtleties, and keenness in doing things. The lion’s face betokens force and strength, and boldness or magnanimity. As the sight of an ox or a calf betokens enduring of labor. The eagle and the six wings signify swiftness. For example, God chose the Assyrians or Babylonians, who should destroy Nineveh. These therefore, as it is in Nahum, the Lord prepared and furnished, that they were swifter than eagles, and the rest as you may read in the first and second chapters of Nahum. And so all creatures are ministers of God’s judgments, coming out of his judicial throne.

\index{Arethas of Caesarea}Then it is also addressed what those beasts do. They go around the throne, always waiting for God’s command, that they may apply it cheerfully, speedily, and stoutly. Neither do they have any rest—mark how he says “have,” not “shall have” or “have had,” but “have”—any rest; that is to say, they are in continual service of God. But let us not understand from this that they are burdened with any painfulness. And also they honor God with continual praise. Arethas: it signifies, he says, no laborious thing. And they have no rest, but a continual occupation, about the singing of godly praises, and so forth.

Finally, here is also presented the hymn and praise of all creatures. In old times David sang, “Praise the sun and moon, \&c.” The same hymn is found in Isaiah 6. And what do all creatures commend in God, whose service God uses, and whose power and operation they feel? Primarily, holiness. These things chiefly concern the substance of the matter. For they teach God to be holy, unspotted, just, good, omnipotent, doing all things, eternal, the beginning of things, and preserver. For they say, “Holy Lord God Almighty, which was, \&c.” Which words we indeed explained in the first chapter. Who would not gather from this that the works and judgments of him are most holy and just? Who therefore shall hereafter reprove the judgments and works of the Lord? Just is the Lord in all his ways, and holy in all his works. This testimony of all creatures makes us willing, ready, cheerful, and untroubled, that we should willingly quiet ourselves in the judgments of God, and murmur at him in nothing, questioning why he should do this or that? But wholly submit ourselves unto God, believing all his works to be good, and to be done for the profit of the godly, and for the most just punishment of the wicked. Holy is God the Father, holy is God the Son, and holy is God the Holy Ghost, holy is one God in Trinity, blessed forevermore. Holy are all his works, and his ways undefiled. And we read more correctly three times holy than nine times, following the example of the Complutensian book. For the former reading, the prophet Isaiah approves. To God Almighty be praise and glory.
